%%%%%%%%%%%%%%%%%%%%%%%%%%%%%%%%%%%%%%%%%
% Beamer Presentation
% LaTeX Template
% Version 1.0 (10/11/12)
%
% This template has been downloaded from:
% http://www.LaTeXTemplates.com
%
% License:
% CC BY-NC-SA 3.0 (http://creativecommons.org/licenses/by-nc-sa/3.0/)
%
%%%%%%%%%%%%%%%%%%%%%%%%%%%%%%%%%%%%%%%%%

%----------------------------------------------------------------------------------------
%	PACKAGES AND THEMES
%----------------------------------------------------------------------------------------

\documentclass[c]{beamer}
%\documentclass[notes]{beamer}
\setbeamertemplate{note page}[show only notes]
\input{../OR_common.tex}
%%%%%%%%%%%%%%%%%%%%%%%%%%%%%%%%%%%%%%%%%%%%%%%%%%%%%%%%%%%%%%%%%%%%%%%%%%%%%
%%%%%%%%%%%%%%%%%%%%%%%%%%%%%%%%%%%%%%%%%%%%%%%%%%%%%%%%%%%%%%%%%%%%%%%%%%%%%
%%%%%%%%%%%%%%%%%%%%%%%%%%%%%%%%%%%%%%%%%%%%%%%%%%%%%%%%%%%%%%%%%%%%%%%%%%%%%

\title[Introduction]{Unit 3. Sensitivity Analysis}

\author{Jordi Villà i Freixa}
\institute[FCTE]{
Universitat de Vic - Universitat Central de Catalunya \\
Study Abroad. Operations Research\\
\medskip
\textit{jordi.villa@uvic.cat}
}
\date{9/05, 2023}
\logo{\includegraphics[width=.1\textwidth]{FCTE}}
\begin{document}

\begin{frame}
\titlepage
\end{frame}


\begin{frame}
    \frametitle{Preliminary}
    This course is strongly based on the monography on Operations Research by Carter, Price and Rabadi \cite{carter}, and in material obtained from different sources (quoted when needed through the slides).
\end{frame}

%%%%%%%%%%%%%%%%%%%%%%%%%%%%%%%%%%%%%%%%%%%%%%%%%%%%%%%%%%%%%%%%%%%%%%%%%%%%%
%%%%%%%%%%%%%%%%%%%%%%%%%%%%%%%%%%%%%%%%%%%%%%%%%%%%%%%%%%%%%%%%%%%%%%%%%%%%%
%%%%%%%%%%%%%%%%%%%%%%%%%%%%%%%%%%%%%%%%%%%%%%%%%%%%%%%%%%%%%%%%%%%%%%%%%%%%%

\begin{frame}
\frametitle{Learning outcomes}
\begin{itemize}
  \item Understanding the concept of postoptimality analysis
  \item Understanding and solving LP problems requiring sensitvity analysis
  \item Understanding the matrix representation of the Simplex solution
\end{itemize}
\end{frame}

\begin{frame}{The concept}
\begin{itemize}
  \item After an optimal solution is found, the analyst needs to review the problem parameters and the solution.
  \item This process is called {\bf postoptimality analysis}:
  \begin{itemize}
    \item confirming or updating problem parameters (cost and availabilities of activities and resources)
    \item if changes need to be introduced in the original parameters, assessing their impact on the optimality of the solution.
  \end{itemize}
  \item If changes are smallre-optimization may not be needed.
  \item {\bf Sensitivity analysis} is the study of the effect that types, ranges , and magnitude of changes in problem parameters have in the value of the objective function, {\em without the need to solving again the new linear problem}.
\end{itemize}
\end{frame}

\begin{frame}{Two types of parameter modifications}
  In a LP problem, I can have two situations
\begin{eqnarray*}
 \text{max }\quad\uvec{c}^t \uvec{x}\\
 \text{subject to }\quad A\uvec{x}\leq \uvec{b}, \forall x_i\geq0
\end{eqnarray*}

I may have interest in knowing what happens with modifications in the parameters $c$ or modifications in the parameters $b$.

\end{frame}

\begin{frame}{Case 1: sensitivity with respect to $c$}

  \begin{Exercise}
    A farmer wants to minimize the cost of the food given of two different types of nutrients $A$ and $B$, provided that she needs a minimum nutrition to be achieved.
    \begin{center}
    \begin{tabular}{c|c|c|c}
     & A & B & price\\
     \hline
     Feed 1 & 10 & 3 & 16\\
     Feed 2 & 4 & 5 & 14\\
     \hline
     requirements & 124 & 60 & \\
    \end{tabular}
    \end{center}

    Find how resilient is she to the changes in the price?

  \end{Exercise}
\note{veure  \url{https://youtu.be/o1pznRt_-y0?t=3378}}
\end{frame}

\begin{frame}{Case 2: sensitivity with respect to $b$}

  In such cases we will explore the coeficients of the objective function in the dual problem, instead.

  \begin{Exercise}
    Given the LP problem
    \begin{equation*}
      \begin{aligned}
        \text{maximize } \quad & z = 16x_1+14x_2 \\
        \text{subject to }\quad &
        \begin{array}{rcl}
          10x_1 + 4x_2 &\geq &124 \\
          3x_1 + 5x_2 &\geq &60 \\
          x_1,x_2 &\geq& 0
        \end{array}
      \end{aligned}
    \end{equation*}

    Explore the sensitivity of the minimum value of $z$ with respect to the parameters in the RHS of the contraints.

  \end{Exercise}

\end{frame}



\note{venim de \url{https://youtu.be/o1pznRt_-y0?t=1603} i \url{https://www.youtube.com/watch?v=yU8updOR87c}}

\begin{frame}{Shadow price is the solution of the dual problem}

  Solving the dual problem gives a lot of insight into the actual situation.

  \begin{Exercise}
  In a company that manufactures two types of cans, $A$ and $B$, the owners want to maximize the production, which dependence on a restricted amount of aluminum, the time the machines can devote to the work and the manpower is given the following table:
  \begin{center}
  \begin{tabular}{c|c|c|c}
   & A & B & limit\\
   \hline
   Aluminium & 50 & 30 & 2000\\
    Machine time & 6 & 5 & 300\\
    Labor & 3 & 5 &200\\
   \hline
   price & 124 & 60 & \\
  \end{tabular}
  \end{center}


  Given an optimal solution of the problem, how the objective function will depend on variations in the different constraints?
  \end{Exercise}

\end{frame}

\begin{frame}{Some tools that help interpreting the calculations}

  Find \href{https://github.com/JordiVillaFreixa/ORcode/blob/main/linear_programming.py}{\em here} a code with the solution using ORtools.

  The Senstivity analysis tool in Excel is nicely explained \href{https://www.youtube.com/watch?v=zKqU5NGE-t0&list=PLjiMsqjDUvBiArYMqCGZSDJNgEMPALdQu&index=1}{\em here}.

\end{frame}


\section{References}
\begin{frame}{References}
    \footnotesize
    \begin{thebibliography}{99}
    \setbeamertemplate{bibliography item}[text]
      \begin{columns}[t]
        \begin{column}{.45\textwidth}
            \bibitem{carter} Michael W. Carter, Camille C. Price, and Ghaith Rabadi. Operations Research, 2nd Edition. CRC Press.
            \bibitem{harel} David Harel, with Yishai Feldman. Algorithmics: the spirit of computing, 3rd Edition. Addison-Wesley.
            \bibitem{rardin} Ronald L. Rardin. Optimization in Operations Research, 2nd Edition. Pearson.
            \bibitem{hefferon} J. Hefferon. \href{http://joshua.smcvt.edu/linearalgebra}{Linear algebra (4th Ed)}.
        \end{column}
        \begin{column}{.45\textwidth}
            \bibitem{riley} K.F. Riley, M.P. Hobson, S.J. Bence. Mathematical Methods for Physics and Engineering (2nd Ed). McGraw Hill.
            \bibitem{nocedal} J. Nocedal, S. J. Wright. Numerical Optimization (2nd Ed). Springer.
            \bibitem{beers} Kenneth J. Beers. Numerical methods for chemical engineering: applications in Matlab. Cambridge University Press.
            \bibitem{barber} D. Barber. Bayesian reasoning and machine learning. Cambridge University Press.
        \end{column}
      \end{columns}
    \end{thebibliography}
\end{frame}
%----------------------------------------------------------------------------------------

\end{document}
